\documentclass[12pt]{article}
\usepackage[utf8]{inputenc}
\usepackage[T1]{fontenc}
\usepackage{hyperref}
\title{Kalispel-Pend d'Oreille Language Resources for K-12 Curriculum Development}
\author{Compiled by Claude Sonnet 5.5 (claude-sonnet-5-5) for Andr\'as Kornai}
\date{}
\begin{document}
\maketitle

\section{Introduction}
Kalispel-Pend d'Oreille (Montana Salish, Salish, Q{\l}ispe; ISO 639-3 fla) is a Salishan language of Montana and Washington. This dossier rests on one search and the resources below are the ones that turned up.

\section{Primary Educational Resources}
The search reported that the language is taught at the Nkwusm Salish Immersion School in Arlee, Montana, and in Kalispell public schools (secondary report only, not verified at a primary source). No school website was found.

\section{Linguistic and Reference Materials}
\begin{itemize}
\item Salish -- Qlipse (Pend d'Oreille) Online Dictionary, UBC knowledgebase entry (developed by Native Teaching Aids): \url{https://knowledgebase.arts.ubc.ca/?p=8454}
\item Same dictionary, UBC wiki page: \url{https://wiki.ubc.ca/Documentation:RelLex/Salish_%E2%80%93_Qlipse_(Pend_d%E2%80%99Oreille)_Online_Dictionary}
\item A Dictionary of the Kalispel or Flat-Head Indian Language (1879, St.\ Ignatius), knowledgebase entry: \url{https://knowledgebase.arts.ubc.ca/?p=3291}
\item Same historical dictionary, UBC wiki page: \url{https://wiki.ubc.ca/Documentation:RelLex/A_Dictionary_of_the_Kalispel_or_Flat-Head_Indian_Language}
\end{itemize}

\section{Religious and Community Resources}
\begin{itemize}
\item Salish-Kalispel materials, ACLS collection (American Philosophical Society guide): \url{https://indigenousguide.amphilsoc.org/entry/39997}
\end{itemize}

\section{Digital and Audio Resources}
The online dictionary is described as multimedia and available as a mobile app. No further audio collections were found.

\section{Curriculum Development Context}
The dictionary draws on Elders and on an earlier Salish language curriculum. The 1879 dictionary is historical and uses older orthography. Contact the tribal language programmes directly for current curriculum materials.
\end{document}